\documentclass[12pt,letterpaper,onecolumn]{article}
\usepackage[OT1]{fontenc}
\usepackage{amsmath,amssymb,amsthm,bm}
\usepackage{microtype}
\usepackage[margin=1in]{geometry}
\usepackage[unicode,bookmarksnumbered,breaklinks,colorlinks=true,
 linkcolor=black,citecolor=black,urlcolor=black]{hyperref}
\newtheorem{theorem}{Theorem}[section]
\newtheorem{lemma}[theorem]{Lemma}
\newtheorem{corollary}[theorem]{Corollary}
\theoremstyle{remark}
\newtheorem{remark}[theorem]{Remark}
\newcommand{\F}{\mathbb F}
\newcommand{\E}{\mathbb E}
\newcommand{\wt}{\operatorname{wt}}
\newcommand{\rkcyc}{r}
\hypersetup{pdftitle={Binary Cyclic Codes with Simultaneously Near-Linear Minimum and Dual Distances},pdfauthor={Yang Shifu},pdfsubject={Mathematical proof note},pdfcreator={LaTeX}}
\pagestyle{plain}
\begin{document}
\noindent{\large\bfseries Mathematical proof note}\par\medskip
\noindent{\raggedright\fontsize{24}{28}\selectfont\bfseries Binary Cyclic Codes with Simultaneously Near-Linear Minimum and Dual Distances\par}\medskip
\noindent{\fontsize{16}{20}\selectfont Yang Shifu}\par\smallskip
\noindent\begin{minipage}{\linewidth}\small National Mobile Communications Research Laboratory and the Frontiers Science Center for Mobile Information Communication and Security, Southeast University, Nanjing 211189, China\end{minipage}\par\medskip
\noindent Version 1.0 --- October 16, 2026\par\bigskip

\section*{Result overview}
For every odd integer $m\ge3$, put $n=2^m-1$. We give a self-contained probabilistic proof that ordinary binary cyclic codes of dimension $(n-1)/2$ can have both minimum distance and dual distance at least $(1/2-o(1))n/(m\log_2m)$. A finite-length sufficient condition is $H_2(\delta)+\delta\le1/(2m)$. The proof counts only low-weight words whose zeroth coordinate is one. Shift symmetry contributes a factor $\delta$ to the count, and a first-moment estimate for random reciprocal-factor choices completes the argument. A probability-tending-to-one version has the same leading distance constant.

\section{Finite-Length Statement}
Identify $\F_2^n$ with $R_n=\F_2[X]/(X^n-1)$ by the coefficient vector. An ordinary binary cyclic code is an ideal of $R_n$. The dual is taken with respect to the usual binary dot product, and $d(A)$ denotes the minimum Hamming weight of a nonzero word of $A$. Write
\[
 H_2(x)=-x\log_2x-(1-x)\log_2(1-x),
 \qquad 0<x<1.
\]

\begin{theorem}\label{thm:finite}
Let $m\ge3$ be odd and $n=2^m-1$. If $0<\delta<1/2$ and
\begin{equation}
 H_2(\delta)+\delta\le\frac1{2m},
 \label{eq:condition}
\end{equation}
then there is an ordinary binary cyclic code $C$ satisfying
\begin{equation}
 \dim C=\frac{n-1}{2},\qquad C\subset C^\perp,\qquad
 \min\{d(C),d(C^\perp)\}>\delta n.
 \label{eq:finite}
\end{equation}
\end{theorem}

Condition~\eqref{eq:condition} is sufficient; no necessity assertion is made. In particular, the distance guarantee is the integer bound $\lfloor\delta n\rfloor+1$.

\section{A First-Coordinate Counting Lemma}
For $f\in R_n$, define its cyclic rank by
\begin{equation}
 \rkcyc(f)=\dim_{\F_2}(R_nf)
 =n-\deg\gcd(f,X^n-1).
 \label{eq:rank}
\end{equation}
For $1\le r\le n$, set
\[
 \mathcal S_r=\{f:\rkcyc(f)=r,\ 0<\wt(f)\le\delta n\},
 \qquad
 \mathcal S_r^{(1)}=\{f\in\mathcal S_r:f_0=1\}.
\]

\begin{lemma}\label{lem:count}
For $0<\delta<1/2$,
\begin{equation}
 |\mathcal S_r^{(1)}|\le\delta\,2^{2rH_2(\delta)}.
 \label{eq:count}
\end{equation}
\end{lemma}
\begin{proof}
The first $k$ coordinates form an information set for every $k$-dimensional binary cyclic code. Indeed, if $g$ is its monic generator, every word has a unique expression $ug$ with $\deg u<k$. Since $g(0)=1$, the first $k$ coefficients of $ug$ determine those of $u$ by a triangular system with diagonal entries one.

If $\rkcyc(f),\rkcyc(h)\le r$, then $f+h$ belongs to the cyclic code $R_nf+R_nh$, of dimension at most $2r$. Hence projection onto the first $\ell=\min(2r,n)$ coordinates is injective on the union of all words of cyclic rank at most $r$.

Assume $\mathcal S_r\ne\varnothing$ and let $F$ be uniform on $\mathcal S_r$. This set is shift invariant, so every coordinate has the same probability
\[
 p=\Pr(F_i=1)=\frac{\E\wt(F)}n\le\delta.
\]
Prefix injectivity and entropy subadditivity give
\[
 \log_2|\mathcal S_r|
 =H(F_0,\ldots,F_{\ell-1})
 \le\ell H_2(p)\le2rH_2(\delta).
\]
Finally, $|\mathcal S_r^{(1)}|=p|\mathcal S_r|$, which proves~\eqref{eq:count}. The empty case is immediate.
\end{proof}

Only $\mathcal S_r$, not $\mathcal S_r^{(1)}$, is asserted to be shift invariant. The argument uses no independence among word coordinates.

\section{Reciprocal Factors and One First Moment}
\begin{proof}[Proof of Theorem~\ref{thm:finite}]
For a polynomial with constant term one, let $P^*(X)=X^{\deg P}P(X^{-1})$ denote its reciprocal. Since $n=2^m-1$ is odd, $X^n-1$ is square-free and splits over $\F_{2^m}$. Its factorization over $\F_2$ has the form
\begin{equation}
 X^n-1=(X+1)\prod_{i=1}^{t}P_iP_i^*,
 \qquad P_i\ne P_i^*,\quad \deg P_i\le m.
 \label{eq:factorization}
\end{equation}
To check the distinct reciprocal pairing, suppose a nontrivial root $\alpha$ and $\alpha^{-1}$ were Frobenius conjugate. If $s>1$ is the order of $\alpha$, then $2^j\equiv-1\pmod s$ for some $j$. But the order of $2$ modulo $s$ divides the odd integer $m$, whereas $-1$ has order two, a contradiction.

Independently and uniformly choose $Q_i\in\{P_i,P_i^*\}$ and put
\[
 g=\prod_iQ_i,\qquad D=\langle g\rangle,
 \qquad C=\langle(X+1)g\rangle.
\]
The dual generator formula yields
\[
 g_{D^\perp}=\left(\frac{X^n-1}{g}\right)^*=(X+1)g.
\]
Thus $C=D^\perp\subset D$, $\dim C=(n-1)/2$, and
\begin{equation}
 \min\{d(C),d(C^\perp)\}=d(D).
 \label{eq:both}
\end{equation}

Fix $f\ne0$. If neither factor in some pair $P_i,P_i^*$ divides $f$, then $\Pr(f\in D)=0$. Otherwise, let $b$ be the number of pairs in which exactly one factor divides $f$. Each such pair forces one independent choice, so $\Pr(f\in D)=2^{-b}$. By~\eqref{eq:rank}, the cyclic rank is the sum of the degrees of the irreducible factors that do not divide $f$. The factor $X+1$ contributes at most one, and therefore $\rkcyc(f)\le1+mb$. In all cases,
\begin{equation}
 \Pr(f\in D)\le2^{-(\rkcyc(f)-1)/m}.
 \label{eq:inclusion}
\end{equation}

Count only the low-weight words with zeroth coordinate one:
\[
 Z_\delta(D)=|\{f\in D:f_0=1,\ \wt(f)\le\delta n\}|.
\]
Any nonzero word can be shifted to such a coordinate position. Hence the exact equivalence is
\begin{equation}
 d(D)>\delta n\quad\Longleftrightarrow\quad Z_\delta(D)=0.
 \label{eq:equivalence}
\end{equation}
Let $a=1/m-2H_2(\delta)$. If $a>0$, Lemma~\ref{lem:count}, \eqref{eq:inclusion}, and linearity of expectation imply
\begin{align}
 \Pr\{d(D)\le\delta n\}
 &\le\E Z_\delta(D)\notag\\
 &\le\sum_{r=1}^{n}\delta\,2^{2rH_2(\delta)}2^{-(r-1)/m}\notag\\
 &\le\frac{\delta\,2^{1/m}}{2^a-1}.
 \label{eq:moment}
\end{align}
Under~\eqref{eq:condition}, $a\ge2\delta$. Since $2^{2\delta}-1>2\delta\ln2$,
\begin{equation}
 \E Z_\delta(D)
 \le\frac{\delta\,2^{1/m}}{2^{2\delta}-1}
 <\frac{2^{1/m}}{2\ln2}
 \le\frac{2^{1/3}}{2\ln2}<1.
 \label{eq:lessone}
\end{equation}
The nonnegative integer $Z_\delta(D)$ is therefore zero for at least one choice. Equations~\eqref{eq:both} and~\eqref{eq:equivalence} finish the proof.
\end{proof}

\section{Asymptotic Distance and Success Probability}
\begin{corollary}\label{cor:asymptotic}
As $m\to\infty$ through odd integers, there are ordinary binary cyclic codes $C_m$ of length $n=2^m-1$ and dimension $(n-1)/2$ such that
\begin{equation}
 \min\{d(C_m),d(C_m^\perp)\}
 \ge\left(\frac12-o(1)\right)\frac{n}{m\log_2m}.
 \label{eq:asymptotic}
\end{equation}
\end{corollary}
\begin{proof}
Let $\delta_m\in(0,1/2)$ be the unique solution of
\[
 H_2(\delta_m)+\delta_m=\frac1{2m}.
\]
The left side is strictly increasing on $(0,1/2)$. Since
\[
 H_2(x)+x=x\log_2(1/x)+O(x)\qquad(x\downarrow0),
\]
its inverse satisfies $\delta_m=(1+o(1))/(2m\log_2m)$. Apply Theorem~\ref{thm:finite}.
\end{proof}

The random choices themselves can be good with probability tending to one at the same leading distance constant. Put
\[
 \ell=\log_2m,\qquad \epsilon_m=\ell^{-1/2},\qquad
 H_2(\widetilde\delta_m)=\frac{1-\epsilon_m}{2m}.
\]
Then $\widetilde\delta_m=(1+o(1))/(2m\ell)$ and $a=\epsilon_m/m$. From~\eqref{eq:moment},
\begin{equation}
 \Pr\{d(D)\le n\widetilde\delta_m\}
 \le\frac{2^{1/m}m\widetilde\delta_m}{(\ln2)\epsilon_m}
 =O(\ell^{-1/2})=o(1).
 \label{eq:whp}
\end{equation}
This version uses~\eqref{eq:moment} directly and does not require~\eqref{eq:condition}.

\section*{Scope of the conclusion}
The equivalence~\eqref{eq:equivalence} is exact. The first-moment condition $\E Z_\delta<1$ and the scalar condition~\eqref{eq:condition} are only sufficient. The result is a probabilistic existence theorem, not a deterministic efficient procedure for choosing a good generator. The relative distance guaranteed by~\eqref{eq:asymptotic} tends to zero; no positive constant relative distance is proved. No orbit-size estimate, rank--weight inequality, nested all-rate construction, self-dual lifting, or channel-capacity theorem is needed for this proof.
\end{document}
